\documentclass[12pt, letterpaper]{article}
\usepackage[margin=1.2in]{geometry}
\usepackage{ebgaramond}
\usepackage[T1]{fontenc}
\usepackage{microtype}
\usepackage{setspace}
\usepackage{parskip}
\usepackage{titling}
\usepackage{fancyhdr}

\setlength{\parindent}{2em}
\onehalfspacing

\pagestyle{fancy}
\fancyhf{}
\fancyhead[L]{\small\textit{The First Picture}}
\fancyhead[R]{\small\textit{NASA Mission Series v1.6}}
\fancyfoot[C]{\small\thepage}
\renewcommand{\headrulewidth}{0.4pt}

\title{\Huge\textbf{The First Picture}\\[0.5em]\large\textit{A short story from Mission 1.6 --- Explorer 6}}
\author{NASA Mission Series\\[0.3em]\small Miles Tiberius Foxglove --- Mukilteo, WA}
\date{March 30, 2026}

\begin{document}
\maketitle
\thispagestyle{empty}

\vspace{1em}
\noindent\rule{\textwidth}{0.4pt}
\vspace{1em}

\noindent The strip-chart recorder had been running all morning, its pen scratching a flat line across the paper with the sound of a mouse running on wood. Flat line. Flat line. The pen moved at a fixed rate, one centimeter per second, and the paper fed through from a roll the size of a man's forearm, and the line it drew was the sound of empty sky.

Jerry Hanson stood in the tracking station at South Point, Hawaii, with his hands in his pockets and his eyes on the recorder. The station was a cinder-block building at the southernmost tip of the Big Island, surrounded by scrub grass and lava rock and wind. The antenna outside --- a Yagi array bolted to a concrete pedestal --- pointed northeast, tracking an object the size of a basketball that was seventeen thousand kilometers above the Indian Ocean and climbing.

Explorer 6. Designation S-2. Launched from Cape Canaveral seven days ago aboard a Thor-Able rocket, carrying instruments that weighed less than a cocker spaniel. The spacecraft was a sphere with four paddle-shaped solar cells, spinning at 2.8 revolutions per minute, and among its instruments was something that had never been sent to orbit before: a photocell behind a narrow slit. A camera. Not a camera as the word was commonly understood --- not a lens focusing light onto film or a tube scanning a phosphor screen --- but a single photosensitive element that measured brightness one point at a time. One revolution of the spacecraft swept the photocell across 360 degrees of whatever was in front of it. One sweep. One line. The brightness at each angular position was converted to a voltage, the voltage was frequency-modulated onto the telemetry carrier, and the carrier was transmitted to the ground.

At the ground, the receiver demodulated the carrier, extracted the brightness data, and fed it to the strip-chart recorder. One line of brightness per revolution. 21.4 seconds per line. The pen moved, the paper advanced, and if the photocell happened to be pointing at something bright, the pen deflected upward. If it was pointing at dark space, the pen stayed flat.

Jerry had been watching the flat line for forty minutes.

\bigskip
\centerline{$\bullet$\quad$\bullet$\quad$\bullet$}
\bigskip

The geometry had to be right. That was the thing people outside the program did not understand. You could not just point a camera at Earth and take a picture. Explorer 6 was spin-stabilized --- it rotated around its own axis for gyroscopic stability, and the spin axis was fixed relative to the stars, not relative to Earth. The photocell's slit swept a great circle in the sky with each revolution, and that great circle might or might not intersect the disk of Earth depending on the spacecraft's attitude, its position in orbit, and the position of the Sun. You needed the sunlit face of Earth to fall within the photocell's scan path at the same time that the spacecraft was within radio range of a ground station. Three conditions, all simultaneously. The orbital mechanics team at Space Technology Laboratories had computed a window: August 14, 1959, between 08:00 and 09:30 local time, when Explorer 6 would be passing over the Pacific with its photocell sweeping across the day-lit crescent of Earth while South Point had a clear line of sight to the spacecraft.

It was 08:47.

The flat line twitched.

Jerry leaned forward. The pen had deflected upward by perhaps two millimeters, held for a fraction of a second, and returned to baseline. He stared at it. The paper continued to feed. The pen scratched. Flat. Flat. Then another twitch --- this one slightly wider, three millimeters of deflection, lasting perhaps a second before it fell back.

``Bill,'' Jerry said, without turning around. ``Come look at this.''

Bill Ostrander was at the receiver rack, adjusting the AGC threshold. The signal from Explorer 6 was weak --- the spacecraft was at twelve thousand kilometers and climbing toward apogee at forty-two thousand --- and the automatic gain control was hunting, trying to hold lock on a carrier that faded and surged with every revolution as the spinning antenna pattern swept past the ground station. Bill walked over and looked at the strip chart.

Another deflection. This one was different: it rose, peaked, and fell in a smooth curve that lasted four seconds. The shape was not random. It had a profile --- a gradual rise, a peak, a gradual fall. Like a hill seen from the side.

``That's the limb,'' Jerry said. He said it quietly, as if saying it too loudly might make it stop.

\bigskip
\centerline{$\bullet$\quad$\bullet$\quad$\bullet$}
\bigskip

The next revolution brought another hill. And the next. Each one was 21.4 seconds after the last --- one per spin --- and each one was slightly different. The peaks shifted in position, their amplitudes varied, and the width of the bright region grew and shrank from line to line. This was not noise. Noise was random. This had structure. This had persistence. This had the slow, deliberate regularity of something real being measured by something spinning.

Jerry pulled a chair over and sat down in front of the recorder. He watched the pen trace line after line, each one a horizontal stripe of brightness data. In his mind, he was stacking them. Line one: a narrow brightness peak on the left side of the scan. Line two: a wider peak, shifted slightly right. Line three: wider still. Line four: the brightest yet, a broad curve that the pen traced in a single smooth arc. Line five: narrowing again. Line six: a faint bump, barely above the noise floor.

He was watching a crescent.

It was crude. The photocell had the resolution of a man squinting through fog. The signal was buried in noise that grew worse with every passing minute as the spacecraft climbed toward apogee and the inverse-square law bled the signal down. But the crescent was there. Six scan lines of brightness data, each one a horizontal slice through a sunlit arc, and when you stacked them in your mind --- when you let your brain do what brains do with partial information, filling in the gaps, recognizing the pattern, seeing the shape behind the static --- what you saw was a planet.

Their planet.

``That's Earth,'' Bill said. He was standing behind Jerry now, both of them watching the pen. ``We're looking at Earth from space.''

Nobody had ever said that sentence before. Nobody in all of human history had looked at a representation of Earth made from above the atmosphere using a machine that was, at that moment, orbiting the planet at 6 kilometers per second. Explorers had climbed mountains and seen the curve of the horizon. Balloons had carried cameras to thirty kilometers and photographed the dark sky above the stratosphere. But no one had ever seen their planet as an object --- as a thing floating in space, bright on one side and dark on the other, finite and round and alone against the black.

Six scan lines. A crescent. A photograph assembled one line at a time by a photocell the size of a pencil eraser, spinning in the dark above the Pacific, whispering its data to an antenna on a lava field at the bottom of Hawaii.

\bigskip
\centerline{$\bullet$\quad$\bullet$\quad$\bullet$}
\bigskip

The seventh scan line was noisier than the sixth. The eighth was worse. By the tenth, the brightness peaks were barely distinguishable from the static floor. Explorer 6 was at nineteen thousand kilometers now, deep in the outer Van Allen belt, and the signal power at South Point had dropped by a factor of six since the first clean scan line. The AGC was wide open, the receiver gain was at maximum, and the pen was drawing a jagged, erratic line that contained as much noise as signal.

Jerry watched two more lines --- the eleventh and twelfth --- and then the crescent was gone. Not because Earth had moved out of the photocell's field of view, but because the signal was too weak to carry the data. The information was still being generated on the spacecraft --- the photocell still spun, still measured brightness, still modulated the carrier --- but the whisper could no longer cross the distance. The message was being spoken, but the listener could no longer hear.

He sat back in his chair and looked at the strip chart. Twelve feet of paper had fed through the recorder. On it, in the pen's scratchy blue ink, were twelve horizontal lines, each one a scan across the day-lit face of Earth. The first few were clean and readable. The middle ones were noisy but usable. The last ones were barely there. Together, they formed something that looked, if you squinted at it with the right kind of eyes, like a smudged crescent --- half a circle of brightness against a flat baseline of dark.

It was the worst photograph ever taken. It was the most important photograph ever taken.

\bigskip
\centerline{$\bullet$\quad$\bullet$\quad$\bullet$}
\bigskip

In the days that followed, the image was processed, enhanced, and published. It appeared in newspapers under headlines that called it ``crude'' and ``blurry'' --- which it was. The photocell had an angular resolution measured in degrees, not arc-minutes. No continents were visible. No oceans, no clouds, no features of any kind beyond the crescent shape itself. A child with a flashlight and a tennis ball could have made a more detailed image.

But Jerry kept thinking about the turkey tail.

Before he'd come to Hawaii, before the tracking stations and the telemetry and the strip-chart recorders, he had grown up in Oregon, in a town surrounded by Douglas-fir forest. He remembered walking the trails behind his house after rain, when the nurse logs were dark with moisture and the fungi bloomed. Turkey tail mushrooms --- \textit{Trametes versicolor}, the many-colored bracket --- grew on fallen logs in concentric rings of brown and tan and cream. Each ring was a period of growth, a record of the fungus expanding outward from its point of attachment. The outermost ring was the newest, the growing edge, thin and white and tender. The inner rings were older, harder, their colors deepened by time and weather. The whole bracket was a map of its own history, written in rings.

The strip chart was the same. Each scan line was a ring. The first line was the inner ring --- the first moment of contact, clear and close. Each subsequent line expanded the image outward, adding one more band to the picture, but each band was noisier than the last, degraded by distance the way the outer rings of a bracket fungus are weathered by exposure. The image built up the way a turkey tail builds up: ring by ring, band by band, the newest growth always at the edge, the whole structure a record of accumulation over time.

Van Gogh would have understood. Vincent, who painted wheat fields and starry nights with a brush so loaded with pigment that the paint stood up from the canvas in ridges you could feel with your fingertips. Who saw beauty in the texture of the medium, not despite the crudeness but because of it. The thick impasto of \textit{Starry Night} was not a failure of technique --- it was the technique. The visible brushstroke was the point. The painting did not hide the process of its own creation; it celebrated it.

Explorer 6's photograph did not hide its process either. The scan lines were visible. The noise was visible. The degradation from line to line was visible. Every limitation of the technology was embedded in the image, and that was what made it beautiful --- not as a picture of Earth (for that, wait a few more years, wait for TIROS, wait for the astronauts and their Hasselblads), but as a picture of the act of seeing. It was a record of a machine trying to see, reaching across distance with a single photocell, assembling an image one line at a time, losing the signal before the picture was complete. The crudeness was the story.

The bracket fungus does not apologize for its rings. The painting does not apologize for its brushstrokes. The first photograph from space does not apologize for its noise. They are honest objects. They show their work.

\bigskip
\centerline{$\bullet$\quad$\bullet$\quad$\bullet$}
\bigskip

Jerry took the strip chart home with him. Not the original --- that went to Space Technology Laboratories in Los Angeles for processing --- but a copy he made by photographing each section of the paper with a Leica on a copy stand. He developed the film in the station darkroom, printed the images on glossy paper, and trimmed them into a strip that he could hold up to the window and see the scan lines backlit by Hawaiian sunlight.

Twelve lines of a planet. Their planet. Seen for the first time by a machine that weighed 64 kilograms, spinning in an orbit that carried it from 237 kilometers above the surface to 42,400 kilometers out, through both Van Allen radiation belts, through the ring current that it was the first spacecraft to discover, and back again every twelve and a half hours. A machine built by people who had never seen what it saw, transmitting what it found to people who had never imagined what they were about to receive.

The strip-chart paper yellowed over the years. The pen ink faded from blue to a pale gray. But the scan lines were still there --- twelve horizontal stripes of varying brightness, the innermost ones clear and definite, the outermost ones ragged with noise, all of them together forming a crescent that was barely recognizable as anything at all unless you knew what you were looking at.

If you knew, you could see the whole Earth in those twelve lines. You could see the sunlit Pacific, the terminator line where day became night, the thin bright arc of atmosphere at the limb. You could see the distance --- the growing noise was the distance, made visible. You could see the effort --- a photocell spinning in vacuum, a ground station listening in Hawaii, a strip chart recording the first image of home.

It was crude. It was imperfect. It was barely legible. It was the first picture anyone had ever taken of the world from beyond it, and no amount of enhancement or technology or subsequent photographs could take that away from twelve scan lines of brightness data recorded on a strip of paper at the southern tip of Hawaii on August 14, 1959, at 08:47 in the morning, when the pen twitched and Jerry Hanson leaned forward and said, quietly, ``Bill, come look at this.''

\vfill

\begin{center}
\small
\rule{0.3\textwidth}{0.4pt}\\[0.5em]
Mission 1.6 --- Explorer 6 (S-2) --- August 7, 1959\\
Organism: \textit{Trametes versicolor} (turkey tail fungus)\\
Bird: Pacific Slope Flycatcher (\textit{Empidonax difficilis})\\
Dedication: Vincent van Gogh (1853--1890)\\[0.3em]
NASA Mission Series --- tibsfox.com
\end{center}

\end{document}
